% Source: https://artofproblemsolving.com/community/c3h2160442_an_original_configuration
% Stable identity: aops-topic-2160442-q01
\begin{center}

\begin{asy}
usepackage("amsmath");
size(6cm);
defaultpen(fontsize(10pt));
pair A=(0,0), B=2*dir(210), C=3.6*dir(-30);
pair Mbc=(B+C)/2;
pair X=extension(A,A+dir(-90),Mbc,Mbc+rotate(90)*(C-B));
pair u=unit(B-A), v=unit(C-A);
pair M=A+dot(X-A,u)*u, Nn=A+dot(X-A,v)*v;
path rm(pair P, pair p, pair q){ real s=0.16; return (P+s*p)--(P+s*p+s*q)--(P+s*q); }
draw(A--B--C--cycle,linewidth(0.9));
draw(B--M,dashed);
draw(A--X,gray(0.5)+dashed);
draw(Mbc--X,gray(0.5)+dashed);
draw(X--M,red); draw(X--Nn,red);
draw(X--B,blue); draw(X--C,blue);
draw(rm(M,-u,unit(X-M))); draw(rm(Nn,-v,unit(X-Nn)));
dot(A); dot(B); dot(C); dot(X); dot(M,red); dot(Nn,red);
label("$A$",A,N); label("$B$",B,NW); label("$C$",C,NE); label("$X$",X,S);
label("$M$",M,W,red); label("$N$",Nn,NE,red);
label("\mbox{}",(min(currentpicture,user=true)+max(currentpicture,user=true))/2,black); // keeps the text after the figure black when LaTeX typesets the labels inline
\end{asy}

\end{center}
\textbf{$N$ lies on ray $AC$, and $AM=AN$.} Since $X$ is on the bisector of $\angle BAC$, the angles $\angle XAB$ and $\angle XAC$ both equal $\dfrac{1}{2}\angle BAC<90^\circ$. In the right triangles $AMX$ and $ANX$ the angle at $A$ is therefore acute, so the feet $M$ and $N$ lie on the rays $AB$ and $AC$ (not on their extensions beyond $A$). These two right triangles have the common hypotenuse $AX$ and equal angles at $A$, so they are congruent:
\[AM=AN\qquad\text{and}\qquad XM=XN .\]

\textbf{$BM=CN$.} Since $X$ is on the perpendicular bisector of $BC$, $XB=XC$. The right triangles $XMB$ and $XNC$ (right angles at $M$ and $N$) have equal hypotenuses $XB=XC$ and equal legs $XM=XN$, so by the hypotenuse-leg criterion they are congruent, and $BM=CN$.

\textbf{Conclusion.} Since $M$ is beyond $B$, $AM=AB+BM$ with $BM>0$. On ray $AC$ the point $N$ is either beyond $C$ or equal to $C$, giving $AN=AC+CN$, or strictly between $A$ and $C$, giving $AN=AC-CN$. In the first case
\[AB+BM=AM=AN=AC+CN=AC+BM,\]
so $AB=AC$, which is excluded. (Indeed, if $AB=AC$ the bisector of $\angle A$ would coincide with the perpendicular bisector of $BC$, and there would be no single point $X$.) Hence $N$ lies strictly between $A$ and $C$. $\blacksquare$

\emph{Remark.} From $AB+BM=AC-CN$ and $BM=CN$ we even get $BM=CN=\dfrac{AC-AB}{2}$, so $M$ beyond $B$ happens exactly when $AB<AC$. The proof does not use the hypothesis $\angle BAC>90^\circ$: in fact $X$ is the midpoint of the arc $BC$ of the circumcircle not containing $A$, which lies outside the triangle for every triangle with $AB\ne AC$.
